\chapter{Genomic Surveillance Graphs}\label{ch:genomic_networks}

\section{Introduction}\label{sec:ch4_introduction}

Chapter~\ref{ch:literature} described how transmission, viral ancestry, and surveillance determine which genomic relationships can be observed. Scottish and UK genomic studies have documented multiple introductions and subsequent circulation of distinct viral lineages \citep{daSilvaFilipe2020,duPlessis2021,Vohringer2021}. Against this background, this chapter applies EpiLink, developed in Chapter~\ref{ch:epilink}, to examine how genetic and temporal compatibility is distributed between sampled population groups and residential areas, and how that structure varies across surveillance periods and lineages.

We assemble 322,366 high-quality sequence records linked to demographic, geographic, deprivation, testing, vaccination, and policy information. The analysis addresses three connected questions. First, who entered the sequenced record, and how did its composition and surveillance conditions change? Second, was compatibility weight concentrated within population or residential categories beyond the reference defined by their shares of total node strength? Third, how did these patterns vary between surveillance windows and viral lineages? Descriptions of graph size and connectivity alongside sequencing volume and coverage provide the sampling context for these comparisons.

The literature motivates exploratory expectations. Regional introductions and circulation make geographic structure plausible, while age-assortative contact patterns motivate examining age structure \citep{daSilvaFilipe2020,duPlessis2021,Mossong2008}. A compatibility graph need not reproduce a contact network, however, and differences in infection risk or ascertainment do not establish an expected direction of sex or deprivation assortativity. These expectations guide interpretation rather than constitute prespecified confirmatory hypotheses.

This chapter systematically characterises EpiLink compatibility graphs across the changing Scottish surveillance record. Demographic and residential attributes add context to graphs constructed from genetic differences and collection dates. The analysis establishes which population patterns are present in the sampled graphs and how consistently they appear across contexts. It informs the interpretation of clusters in Chapters~\ref{ch:cluster_transition_screening} and~\ref{ch:candidate_characterisation}, while leaving individual transmission links and exposure settings to be assessed with independent epidemiological evidence.

\section{Materials and Methods}\label{sec:ch4_methods}

We combined seven data sources into a single clustering dataset that underpins the analyses here and in Chapters~\ref{ch:cluster_transition_screening} and~\ref{ch:candidate_characterisation}. For each genomic sequence, it provides a cluster assignment together with linked demographic, clinical, geographic, and deprivation attributes, alongside epidemic and public-health policy variables describing the surveillance context at the time of sample collection.

We analysed sequences within rolling 21-day surveillance windows, advanced in seven-day steps (Figure~\ref{fig:rolling_window_schematic}). The 21-day duration accommodates multiple SARS-CoV-2 transmission generations and allows for infection-to-detection delays. The weekly step aligns with reporting cycles, while the 14-day overlap reduces the effect of where window boundaries fall. We analysed each Pango lineage separately within each window to avoid uninformative comparisons between lineages with large genetic differences. The rationale for these window and lineage choices is discussed in Chapter~\ref{ch:literature}, drawing on studies of transmission timing and viral kinetics \citep{Alene2021,Hart2021,Xu2023,Puhach2022,OordtSpeets2024}.

\begin{figure}[htbp]
	\centering
	\includegraphics[width=0.8\textwidth]{rolling_window_schematic}
	\caption[Rolling surveillance window construction]{
		\textbf{Rolling surveillance window construction.} We defined window membership using half-open intervals, which include the start date but exclude the end date.
	}\label{fig:rolling_window_schematic}
\end{figure}

The remainder of this section describes how we link and prepare the seven data sources (Section~\ref{sec:ch4_data}) and construct compatibility graphs and identify clusters within each window and lineage (Section~\ref{sec:ch4_pipeline}). We then explain how we examine whether compatibility links concentrate within shared population categories, how this pattern varies across windows and lineages, and the size and connectivity of the graphs (Section~\ref{sec:ch4_graph_characterisation}). Supplementary methods are provided in Appendix~\ref{sec:app_ch4_supp_methods}.

\begin{tcolorbox}[
    title=\textbf{Glossary}: key terms in graph-based genomic surveillance used in this chapter, 
    halign=right, 
    fontupper=\small
    ]
\begin{description}
    \item[Unique sequence:] One consensus-genome record from a sequenced sample, summarising its SARS-CoV-2 genetic material. Different records can contain identical genomes, and a patient can contribute more than one record.
    \item[Compatibility graph:] A mathematical structure representing pairwise relationships between nodes, here constructed from compatibility scores between unique sequences.
    \item[Nodes:] Each node represents one sequence within a compatibility graph.
    \item[Edges:] Connections between pairs of nodes, weighted by their compatibility score.
    \item[Node degree:] The number of edges connected to a node.
    \item[Node strength:] The sum of edge weights (compatibility scores) connected to a node.
    \item[Complete graph:] A graph where every node is connected to every other node.
    \item[Sparsification:] The process of removing weak edges to simplify the graph while retaining most of the total compatibility weight.
    \item[Communities or clusters:] Groups of nodes identified by the Leiden algorithm as having denser internal connections than external connections.
    \item[Singleton cluster:] A cluster containing one sequence. This occurs when a window-lineage group contains only one sequence, when a sequence has no retained edges, or when Leiden assigns it to a community of one.
\end{description}
\end{tcolorbox}

\subsection{Data sources and preprocessing}\label{sec:ch4_data}

Of the seven data sources (Table~\ref{tab:data_sources}), Datasets II--IV are confidential administrative resources provided by Public Health Scotland (PHS)\nomenclature{PHS}{Public Health Scotland} via the electronic Data Research and Innovation Service (eDRIS)\nomenclature{eDRIS}{electronic Data Research and Innovation Service}. The remaining four datasets (I, V--VII) are publicly available. Dataset~II was assembled by \citet{Gamza2024}; it supplied the sequence, specimen, and patient identifiers needed to connect sequence records to testing and vaccination records. We retained only records linked to high-quality sequences. We applied this quality filter before clustering (Section~\ref{sec:ch4_pipeline}) so that all subsequent steps used the same set of high-quality sequences.

\begin{table}[htbp]
	\centering
	\caption[Seven Scottish data sources]{Seven Scottish data sources used to construct the analytical products here and in Chapters~\ref{ch:cluster_transition_screening} and~\ref{ch:candidate_characterisation}. Datasets I--VII are the raw inputs; processed intermediate files derived from them are documented in the project repository (see the \hyperref[sec:software]{Software and Reproducibility} statement). Access levels are denoted as \textit{Confidential} (PHS/eDRIS) or \textit{Public}; public-source web addresses and derivation notes are given in the source column and table notes.}\label{tab:data_sources}
	\begin{thesistablebody}[\footnotesize\setlength{\tabcolsep}{3pt}\renewcommand{\arraystretch}{1.04}]{@{}P{0.045\linewidth}P{0.175\linewidth}P{0.545\linewidth}P{0.175\linewidth}@{}}
	\toprule
	\textbf{No.} & \textbf{Dataset} & \textbf{Content} & \textbf{Source} \\
	\midrule
	I & COG-UK SARS-CoV-2 consensus sequences
		& A multi-FASTA containing 3,392,463 SARS-CoV-2 consensus genomic sequences, retrieved from the COG-UK archive on 25 October 2023. The FASTA headers contain the COG-UK sequence name which contains the country of origin, allowing filtering to 386,132 Scottish sequences.
		& Public; \textit{CLIMB-COVID}\nomenclature{CLIMB}{Cloud Infrastructure for Microbial Bioinformatics}\textsuperscript{a}. \\
	\addlinespace[0.25em]
	II & PHS sequence metadata table
		& One record per specimen testing positive by polymerase chain reaction (PCR)\nomenclature{PCR}{polymerase chain reaction} with a released genome sequence. Fields used in this analysis included sequence, specimen and patient identifiers, sample collection date, sex, five-year age band, test result, test type (PCR or lateral-flow device; LFD\nomenclature{LFD}{lateral-flow device}), test reason, vaccination date, dose number, booster indicator, and 2011 Data Zone of residence. Raw file: 369,580 records spanning 04 July 2020 to 13 February 2023.
		& Confidential; PHS via eDRIS. \\ 
	\addlinespace[0.25em]
	III & PHS COVID-19 testing records
		& Individual test records from the same source used by \citet{Gamza2024}. It was used here for constructing daily Data Zone-level testing summaries at the date of sample collection.
		& Confidential; PHS via eDRIS. \\ 
	\addlinespace[0.25em]
	IV & PHS COVID-19 vaccination records
		& Individual vaccination event records from the same source used by \citet{Gamza2024}. It was used here for constructing daily Data Zone-level vaccination summaries at each sample collection date.
		& Confidential; PHS via eDRIS. \\ 
	\addlinespace[0.25em]
	V & Scottish Index of Multiple Deprivation 2020 version 2
		& Socioeconomic deprivation attributes for Scotland's 6,976 2011 Data Zones. Fields used in this analysis included total population, overall SIMD rank, and administrative identifiers for urban/rural class, local authority, and National Health Service (NHS)\nomenclature{NHS}{National Health Service} health board.
		& Public; \textit{Scottish Government}\textsuperscript{b}. \\
	\addlinespace[0.25em]
	VI & 2011 Scottish Data Zone boundaries
		& The 2011 Data Zone boundary file (shapefile format), providing polygon geometry, standard area in km\textsuperscript{2}, and geographic centre points weighted by population distribution in the British National Grid (OSGB36\nomenclature{OSGB36}{Ordnance Survey Great Britain 1936}, EPSG\nomenclature{EPSG}{European Petroleum Survey Group}:27700).
		& Public; \textit{Scottish Government}\textsuperscript{c}. \\
	\addlinespace[0.25em]
	VII & Oxford COVID-19 Government Response Tracker national indices (OxCGRT)\nomenclature{OxCGRT}{Oxford COVID-19 Government Response Tracker}
		& Tables providing daily Scotland-level values for the Stringency Index and the Containment and Health Index.
		& Public; \textit{OxCGRT GitHub repository}\textsuperscript{d}. \\
	\bottomrule
	\end{thesistablebody}
	\par\vspace{0.4\baselineskip}
	\begin{minipage}{\linewidth}
	\footnotesize\raggedright %chktex 1
	\textit{Table notes.} Public source web addresses, all accessed 20 August 2026:
	\textsuperscript{a}\url{https://data.covid19.climb.ac.uk/};
	\textsuperscript{b}\url{https://www.gov.scot/collections/scottish-index-of-multiple-deprivation-2020/};
	\textsuperscript{c}\url{https://ckan.publishing.service.gov.uk/dataset/data-zone-boundaries-2011};
	\textsuperscript{d}\url{https://github.com/OxCGRT/covid-policy-dataset/tree/main/data/timeseries_indices}.
	\end{minipage}
\end{table}
\FloatBarrier%

We linked and prepared these datasets (Table~\ref{tab:data_sources}) for the analyses as follows:

\begin{description}
    \item[\textbf{Dataset~I (COG-UK consensus sequences).}] We selected the 386,132 Scottish sequences from the 3,392,463 COG-UK consensus genomes by country of origin. We processed these sequences with Nextclade \citep{Aksamentov2021}, which aligned each sequence to the Wuhan-Hu-1 reference genome (GenBank accession MN908947.3). Nextclade also assigned variant labels (World Health Organization (WHO)\nomenclature{WHO}{World Health Organization} variants of concern where applicable, Nextclade clades, and Pango lineages) alongside quality-control metrics and additional metadata. We retained only sequences with a `good' quality-control status, yielding 322,366 high-quality sequences corresponding to 315,325 distinct patients. This step produced two datasets: \textbf{Dataset~Ia}, consisting of the multi-FASTA alignment of these high-quality sequences, which was used for pairwise genetic distance calculations; and \textbf{Dataset~Ib}, comprising the accompanying table of sequence-level annotations.

    \item[\textbf{Dataset~II (PHS sequence metadata).}] We joined the Nextclade annotation table Dataset~Ib to the PHS sequence metadata to produce \textbf{Dataset~IIa}. All the high-quality sequences had matching PHS metadata. A total of 347 patients (0.11\%) had more than one sequence recorded on the same day. Following \citet{Gamza2024}, we retained these multiple sequence records. Finally, we recoded the 36 original `test reason' entries into 11 broader analytical categories to facilitate analysis (Appendix~Table~\ref{tab:app_test_reason_map}).

	\item[\textbf{Datasets~III and~IV (testing and vaccination records).}] We generated \textbf{Dataset~IIb} by compiling daily and cumulative testing and vaccination summaries from Datasets~III and~IV. These summaries were aggregated by date, Data Zone, surveillance window, or policy period. We then merged them with Dataset~IIa to add local testing and vaccination measures to each sequence record.

	\item[\textbf{Dataset~V (SIMD).}] Following PHS practice, we sorted the 6,976 Data Zones by their overall SIMD rank. We then assigned population-weighted deprivation quintiles to account for differences in Data Zone populations (Quintile 1: most deprived; Quintile 5: least). Population weighting produced quintiles containing approximately one-fifth of the population, with a maximum deviation from 20\% of 0.012 percentage points, compared with 0.706 when grouping equal numbers of Data Zones (Appendix~Figure~\ref{fig:app_simd_population_weighting} and Table~\ref{tab:simd_population_weighting}). Deprivation quintiles were then mapped to each record in Dataset~IIb based on patient residence to produce \textbf{Dataset~IIc}.

	\item[\textbf{Dataset~VI (Data Zone boundaries).}] We added location and area data to Dataset~IIc for the spatial summaries. For each record, we attached the coordinates of the Data Zone's population-weighted centre (EPSG:27700) and its standard area in km\textsuperscript{2}, yielding \textbf{Dataset~IId}.

	\item[\textbf{Dataset~VII (OxCGRT Indices).}] We linked daily Scottish Stringency and Containment and Health Index values \citep{Hale2021} to thirteen policy periods (Appendix~Table~\ref{tab:scot_covid_policy}) curated from outbreak response timelines published by the Scottish Parliament Information Centre (SPICe)\nomenclature{SPICe}{Scottish Parliament Information Centre} \citep{ScottishParliamentSPICe2023} and Scottish COVID-19 Inquiry (SCI)\nomenclature{SCI}{Scottish COVID-19 Inquiry} \citep{ScottishCovidInquiryTimeline}. The two indices were highly correlated in Scotland, so we used the mean Stringency Index as the primary summary of restriction and containment levels (Appendix~Figure~\ref{fig:app_policy_index_comparison}). Policy annotations were then attached to each record in Dataset~IId by collection date to produce the final consolidated sequence metadata table, \textbf{Dataset~IIe}, which contains one record per unique sequence with linked demographic, clinical, geographic, deprivation, vaccination, and policy attributes.
\end{description}

\subsection{Clustering pipeline}\label{sec:ch4_pipeline}

Within each window-lineage subset of Datasets~Ia and~IIe, pairwise Tamura-Nei 1993 (TN93) genetic distances were calculated using the settings described in Section~\ref{sec:ch3_epilink_boston} \citep{TamuraNei1993,KosakovskyPond2018}. To express TN93 distances as approximate numbers of nucleotide differences, we multiplied these proportional distances by the SARS-CoV-2 alignment length (29,903 nucleotides). Pairwise temporal distances represented the difference in collection dates. These genetic and temporal distances were input into the stochastic-genetics EpiLink formulation (ES; Section~\ref{sec:ch3_epilink_variants}), using the default SARS-CoV-2 parameters and recent-linkage target described in Sections~\ref{sec:ch3_parameters_sensitivity} and~\ref{sec:ch3_epilink_model}, respectively. In total, 2,647,110,913 pairwise comparisons were analysed across all window-lineage strata (134 surveillance windows and 790 Pango lineages). Table~\ref{tab:pipeline_summary} summarises the clustering design applied to the Scottish dataset and analysis counts. The resulting graphs and cluster assignments provided the common basis for the analyses in this chapter and Chapters~\ref{ch:cluster_transition_screening} and~\ref{ch:candidate_characterisation}.

To reduce the computational cost of graphs containing every possible edge, we removed edges with very low compatibility scores. We retained only edges with weights exceeding $1\times10^{-3}$; this sparsification threshold preserved 99.0\% of the total compatibility weight while removing 67\% of the edges across all window-lineage strata (Appendix~Figure~\ref{fig:app_parameter_sensitivity}). Retaining most of the total weight does not establish that cluster assignments are unchanged; assessing that effect would require a direct comparison across sparsification thresholds. Using the retained edges, we constructed a weighted, undirected compatibility graph for each window and lineage and partitioned it across a range of Leiden resolution parameters (Table~\ref{tab:pipeline_summary}). The primary analysis used a resolution of $\gamma = 0.3$, selected using the earlier version of the synthetic resolution sweep. Correcting the reference memberships subsequently selected $\gamma = 0.2$ for Boston (Section~\ref{sec:ch3_epilink_boston}); the Scottish analyses reported here retain the original $\gamma = 0.3$ specification. We retained the full parameter range for sensitivity analyses (Appendix~Figure~\ref{fig:app_parameter_sensitivity}).

\begin{table}[htbp]
	\centering
	\caption[Scotland clustering design and analysis counts]{
		Scotland clustering design and analysis counts. \emph{Panel A} summarises the choices governing graph construction and clustering. \emph{Panel B} reports the analysis scope, pairwise comparisons, and cluster counts at the primary Leiden resolution ($\gamma=0.3$).}\label{tab:pipeline_summary}
	\begin{thesistablebody}{@{}P{0.36\linewidth}P{0.61\linewidth}@{}}
		\toprule
		\multicolumn{2}{@{}l}{\textbf{Panel A. Clustering design}} \\
		\midrule
		\textbf{Choice} & \textbf{Specification} \\
		\midrule
		Rolling windows & 21-day width; 7-day step; 14-day overlap. \\
		\addlinespace[0.3em]
		Stratification & Separate compatibility graphs within each rolling window and Pango lineage. \\
		\addlinespace[0.3em]
		Edge retention & Compatibility score $>1\times10^{-3}$; retains 33.1\% of scored comparisons and 99.0\% of total compatibility weight. \\
		\addlinespace[0.3em]
		Community detection & Leiden algorithm; primary $\gamma=0.3$; sensitivity grid $\gamma\in\{0.1,0.2,\ldots,0.8\}$. \\
		\addlinespace[0.3em]
		Singleton handling & Retain single-sequence strata, nodes with no retained edges, and single-node Leiden communities as clusters of size one. \\
		\addlinespace[0.6em]
		\toprule
		\multicolumn{2}{@{}l}{\textbf{Panel B. Analysis scope and outputs}} \\
		\midrule
		\textbf{Summary} & \multicolumn{1}{r}{\textbf{Count}} \\
		\midrule
		Unique high-quality sequence records & \multicolumn{1}{r}{322,366} \\
		Rolling windows & \multicolumn{1}{r}{134} \\
		Pango lineages & \multicolumn{1}{r}{790} \\
		\addlinespace[0.4em]
		Scored pairwise comparisons & \multicolumn{1}{r}{2,647,110,913} \\
		Retained graph edges & \multicolumn{1}{r}{876,277,547} \\
		\midrule
		\multicolumn{2}{@{}l}{\emph{Cluster instances at the primary resolution ($\gamma=0.3$)}} \\
		\midrule
		All clusters & \multicolumn{1}{r}{551,385} \\
		\quad Singletons & \multicolumn{1}{r}{454,907} \\
		\quad Non-singletons & \multicolumn{1}{r}{96,478} \\
		\bottomrule
	\end{thesistablebody}
	\par\vspace{0.4\baselineskip}
	\begin{minipage}{\linewidth}
	\footnotesize\raggedright
	\textit{Table notes.} Pairwise comparisons, retained edges, and cluster instances are counted within window-lineage strata and summed across the analysis. Overlapping windows can therefore contribute the same sequence or sequence pair more than once; the sequence count refers to unique records.
	\end{minipage}
\end{table}

We retained clusters containing one sequence, whether they arose from single-sequence window-lineage groups, nodes with no retained edges, or single-node Leiden communities. The resulting clusters were identified separately within each window and lineage and form the basis for the analyses presented in Chapters~\ref{ch:cluster_transition_screening} and~\ref{ch:candidate_characterisation}.

In the final step, we joined the cluster assignments to Dataset~IIe to produce a single clustering dataset containing one record per unique sequence for each window and Leiden resolution.
Testing reason was summarised descriptively after recoding the 36 raw values into 11 analytical categories. Because missingness and recorded reasons for testing changed substantially over time, testing reason was not used as a predictor in later analyses. Vaccination status was likewise treated as surveillance context. Vaccination dose groups and time since the last recorded dose were described separately in each rolling window.

\subsection{Characterisation of compatibility graphs}\label{sec:ch4_graph_characterisation}

After describing the surveillance background, we examined how compatibility weight was distributed across population categories within graphs for each window and lineage. These graph summaries describe structure among sampled infections; they do not estimate the composition of the underlying infected population or identify realised transmission events.

\subsubsection*{Categorical assortativity coefficients}\label{sec:ch4_assortativity_matrices}

For each window-lineage compatibility graph, we asked whether retained compatibility weight was concentrated within shared population categories. We compared the observed pattern with a random-pairing baseline that preserves each category's share of total node strength, calculated using edges with recorded attribute values at both ends. This asks whether compatibility weight connects sequences in the same category more often than expected from these weight shares, rather than simply describing which categories were most common in the sequenced cohort.

We examined seven categorical attributes: sex; age band (five-year intervals: 00--04, 05--09, \ldots, 70--74, 75+); broader age group (00--04, 05--14, 15--24, 25--64, 65--74, 75+); population-weighted SIMD quintile; urban/rural class; local authority; and health board. We retained both sets of age groups because assortativity values can depend on how categories are defined and how many there are.

For each graph $g$ and attribute $m$, we used only retained undirected compatibility edges for which both sequences had recorded values for that attribute. Each edge contributed in proportion to its compatibility weight. We then formed a weighted mixing matrix $E^{(g,m)}$, counting each undirected edge in both directions and normalising the entries to sum to one. The entry $E^{(g,m)}_{ab}$ describes the share of this weight joining category $a$ to category $b$.

We used Newman's weighted nominal assortativity coefficient to quantify categorical mixing in each compatibility graph \citep{Newman2003}:

\[
	r_{gm}
	=
	\frac{
		\operatorname{tr}\!\left(E^{(g,m)}\right)
		-
		\sum_a \left(p^{(g,m)}_a\right)^2
	}{
		1 -
		\sum_a \left(p^{(g,m)}_a\right)^2
	},
	\qquad
	p^{(g,m)}_a=\sum_b E^{(g,m)}_{ab}.
\]

Here, $p^{(g,m)}_a$ is category $a$'s share of total node strength among edges with recorded values for attribute $m$ at both ends. The trace $\operatorname{tr}\!\left(E^{(g,m)}\right)$, the sum of the diagonal entries, gives the observed share of compatibility weight connecting sequences in the same category. The numerator compares this with the expected share, $\sum_a \left(p^{(g,m)}_a\right)^2$, if edge endpoints were paired at random while preserving these category weight shares. The denominator rescales the difference: positive values indicate more within-category compatibility weight than expected, values near zero indicate little departure from this baseline, and negative values indicate less within-category weight than expected. We treated estimates as undefined when fewer than two categories were observed or the denominator was zero.

We quantified graph-specific estimation uncertainty with an edge-weight multiplier bootstrap, a resampling method that tests how sensitive results are to changes in edge weights \citep{Efron1979,Praestgaard1993}. For each of 500 replicates, we randomly varied the retained edge weights, recomputed the weighted mixing matrix, and recalculated $r_{gm}$. Graph-level standard errors were the standard deviations of the bootstrap estimates that could be calculated, and confidence intervals used percentile intervals. These intervals measure sensitivity to changes in retained edge weights. They do not describe the range expected under random mixing or uncertainty about unsequenced infections. The bootstrap does not remove whole nodes or account for dependence caused by shared endpoints and overlapping rolling windows.

\subsubsection*{Pooling across lineages and windows}\label{sec:ch4_assortativity_meta}

To account for lineage-size differences and temporal overlap between rolling windows, we summarised assortativity in two stages. First, within each rolling window, we pooled lineage-specific estimates for each attribute using random-effects meta-analysis. This allows for variation between lineages. Second, we estimated the overall attribute-level mean across rolling windows using a regression model that accounts for calendar overlap between windows. These pooled means are the primary summaries reported in the results. Full technical details of the meta-analysis and regression models are provided in Appendix~\ref{sec:app_ch4_assortativity_methods}.

\subsubsection*{Variance decomposition}\label{sec:ch4_assortativity_variance_decomp}

The pooled assortativity summaries indicate the average tendency for compatibility weight to connect infections in the same population category, but not whether that pattern is stable across the sampled epidemic. We therefore used variance decomposition to assess how much variation in assortativity was associated with differences between sampling windows, between lineages, or both, and how much remained unexplained. This helps describe how geographic or demographic patterns in compatibility vary over time and between circulating lineages. Full technical details are provided in Appendix~\ref{sec:app_ch4_assortativity_methods}.

\subsubsection*{Node degree and strength assortativity}\label{sec:ch4_assortativity_degree_strength}

Categorical assortativity describes how compatibility weight is distributed across observed population attributes. Node degree and strength assortativity instead describe patterns in the connections themselves (see the Glossary above). We used these additional measures of graph structure to assess whether sequences with many connections or large total connection weights tended to connect to sequences with similarly many connections or similarly large total connection weights. We reported three diagnostics: \emph{unweighted} degree assortativity, \emph{weighted} degree assortativity, and \emph{weighted} strength assortativity. For each rolling window, we computed these diagnostics within lineage-specific graphs and summarised them across lineages (Appendix~\ref{sec:app_ch4_assortativity_methods}).

\section{Results}\label{sec:ch4_results}

The Scottish surveillance dataset contained 322,366 high-quality SARS-CoV-2 sequence records from 315,325 distinct patients. Samples were collected between 4 July 2020 and 10 February 2023. These counts define the sequences available to the graph analyses, not the full history of infection in Scotland. We present the results in two parts. First, we describe the surveillance context: sequencing volume and coverage across policy periods and circulating variants, reasons for testing, and the demographic, geographic, deprivation, and vaccination composition of the sequenced cohort. Second, we characterise the resulting compatibility graphs by population attribute assortativity.

\subsection{Surveillance context and cohort composition}\label{sec:ch4_surveillance_setting}

The sequenced record was generated under changing policy, epidemic, testing, and vaccination conditions. We summarise these changing surveillance conditions across 13 policy periods defined from the published timelines, grouped into six broader epidemic eras, to provide context for interpreting compatibility-graph structure.

\subsubsection*{Sequencing volume and coverage}\label{sec:ch4_volume_coverage}

The earliest samples included in the analysis were collected on 4 July 2020, during Route Map Phase~2 (P2), after the first epidemic wave and during easing of the first national lockdown. Most of the analysis period fell within Route Map Phase~3 (P3) and later periods, when restrictions, testing behaviour, sequencing priorities, vaccination coverage, and viral lineages changed substantially. Figure~\ref{fig:surveillance_timeline} and Table~\ref{tab:policy_period_info} place sequence volume, weekly sequencing coverage, policy periods and epidemic-era stringency, and clade turnover on a common calendar axis.

\begin{figure}[htbp]
    \centering
    \includegraphics[width=\textwidth]{assets/scotland/policy_sequences_over_time}
    \caption[Scottish SARS-CoV-2 surveillance context]{\textbf{Scottish SARS-CoV-2 surveillance context.} (A) Daily sequence volume (bars), 7-day rolling mean (solid line), and weekly sequencing coverage (dashed red line), shown together with the policy-period context. The strip above the panel marks the policy periods, coloured by mean daily OxCGRT Stringency Index; policy-period codes are defined in Table~\ref{tab:policy_period_info}. The broader epidemic-era shading reflects average stringency across the grouped periods, with era labels shown vertically. (B) Weekly composition of all 30 observed Nextstrain clade assignments, including recombinant sequences, over the same period. Colours identify clades; within broader variant groupings, colour intensity reflects total sequence count, with higher-count clades shown darker. The legend is ordered by each clade group's first observed collection date. The thin stacked-bar inset summarises the overall proportion of sequences belonging to each clade across the full dataset (Dataset~IIe from Section~\ref{sec:ch4_data}).
	}\label{fig:surveillance_timeline}
\end{figure}

\begin{table}[htbp]
    \centering
    \caption[Key policies and restriction levels in Scotland during the COVID-19 pandemic]{
    Key policies and restriction levels in Scotland during the COVID-19 pandemic.
    Dates are inclusive. Mean daily OxCGRT Stringency Index values (with standard deviations) summarise the relative restrictiveness of each period, with higher values corresponding to more stringent measures; policy descriptions for these periods are provided in Appendix~Table~\ref{tab:scot_covid_policy}. The epidemic-era column groups adjacent policy periods marked above Figure~\ref{fig:surveillance_timeline}A into the broader phases. Policy periods are listed to 5 May 2023 to provide continuous context; the genomic analyses are restricted to sequences collected up to 10 February 2023. Policy periods ending before 4 July 2020 are omitted because the cohort contains no sequences collected during those periods.
    }\label{tab:policy_period_info}
    \begin{thesistablebody}{@{}llllc@{}}
    \toprule
    \textbf{Epidemic era} & \textbf{Policy period} & \textbf{Period code} & \textbf{Dates} & \textbf{Mean stringency} \\
    \midrule
        Early restriction easing & Route Map Phase 2 & P2 & 19 Jun --- 09 Jul 2020 & $ 76 \pm 2 $ \\
        Early restriction easing & Route Map Phase 3 & P3 & 10 Jul --- 01 Oct 2020 & $ 68 \pm 3 $ \\
        Autumn/winter restrictions & Pre-tier tightening & T1 & 02 Oct --- 01 Nov 2020 & $ 65 \pm 0 $ \\
        Autumn/winter restrictions & Five-tier framework & F5 & 02 Nov 2020 --- 04 Jan 2021 & $ 71 \pm 6 $ \\
        Autumn/winter restrictions & Second lockdown & L2 & 05 Jan --- 01 Apr 2021 & $ 85 \pm 3 $ \\
        Spring/summer 2021 easing  & Stay local (Level 3) & SL & 02 Apr --- 25 Apr 2021 & $ 70 \pm 1 $ \\
        Spring/summer 2021 easing  & Level 3 & L3 & 26 Apr --- 16 May 2021 & $ 58 \pm 0 $ \\
        Spring/summer 2021 easing  & Level 2 / Level 1 & L21 & 17 May --- 18 Jul 2021      & $ 56 \pm 5 $ \\
        Spring/summer 2021 easing  & Level 0  & L0   & 19 Jul --- 08 Aug 2021 & $ 53 \pm 0 $ \\
        Near-normal / Delta & Near-normal & NN & 09 Aug --- 28 Nov 2021  & $ 31 \pm 3 $ \\
        Omicron response & Omicron wave & OM & 29 Nov 2021 --- 23 Jan 2022 & $ 35 \pm 4 $ \\
        Omicron response & Final easing & FE & 24 Jan --- 17 Apr 2022 & $ 20 \pm 7 $ \\
        Post-restriction & Post-restriction & PR & 18 Apr 2022 --- 05 May 2023 & $ 8 \pm 3 $  \\
        \bottomrule
    \end{thesistablebody}
\end{table}

Sequencing volume and coverage varied markedly, but not in parallel, across these periods (Table~\ref{tab:policy_denominators}). Median sequencing coverage exceeded half of confirmed positives during the short spring 2021 easing periods SL (54.7\%) and L3 (53.8\%), but was lower during higher-volume Delta and Omicron periods. The near-normal/Delta period NN had a median of 12,088 genomes per window and 13.4\% coverage, while the final-easing Omicron period FE had the largest median genome count per window (32,026 sequences) but 11.8\% coverage. After population-scale testing was withdrawn, the post-restriction period PR had lower median sequence volume (1,344 genomes per window) and 8.5\% coverage. High sequencing volume therefore did not necessarily imply a high ratio of sequenced genomes to recorded positive tests.

\input{assets/scotland/tab_policy_denominators}

The policy-period strip in Figure~\ref{fig:surveillance_timeline}A uses the OxCGRT Stringency Index to summarise the severity of official restrictions on a 0--100 scale; it does not measure adherence, implementation, effectiveness, or individual exposure. Changes in circulating clades overlapped with this changing policy context, from 20B/20E circulation during early easing, through Alpha and Delta replacement, to rapid changes in the circulating Omicron sublineages. Although graph construction uses Pango lineage stratification, higher-level clades are displayed here for readability. Policy-period and clade labels describe the epidemic context; their separate effects cannot be established here.

\subsubsection*{Reasons for testing}\label{sec:ch4_ascertainment}

The recorded reasons for testing also changed over the study period (Figure~\ref{fig:test_reason_by_policy_era}).

\begin{figure}[htbp]
	\centering
	\includegraphics[width=\textwidth]{fig_test_reason_by_policy_era}
	\caption[Test-reason composition by epidemic era]{
		\textbf{Reason for testing among sequenced infections, by epidemic era.} Each cell gives the number of unique sequences recorded under a given testing reason within an epidemic era; cell shading gives the share of sequences within each epidemic era (column proportion), so that colour emphasises changes in the recorded reasons for testing while the printed values give absolute counts. The 36 raw testing-reason values were recoded into 11 analytical categories (Appendix~Table~\ref{tab:app_test_reason_map}).
		}\label{fig:test_reason_by_policy_era}
\end{figure}

Across all epidemic eras, symptomatic community testing was the dominant recorded reason: ``Symptomatic Citizen'' was the most frequent testing reason in every era except the post-restriction era, peaking at 50,410 sequences during the Omicron response. Confirmatory testing also contributed large volumes during the Omicron response (41,170 sequences), while contact tracing declined to 233 sequences in the post-restriction era. The post-restriction era also contained a large ``Isolation Scheme'' contribution (7,974 sequences) and the highest number of records with no recorded testing reason (17,949 ``Missing''). Given these changes over time in missing testing reasons and in how infections were detected, testing reason is used only descriptively here and is not used as a predictor in later analyses.

\subsubsection*{Cohort composition}\label{sec:ch4_cohort_composition}

Figure~\ref{fig:sequence_composition_by_policy} and Table~\ref{tab:sequence_composition_by_policy} summarise who entered the sequenced record, showing how demographic, geographic, and deprivation categories were distributed across the six epidemic eras.

Females contributed slightly more sequences (174,146; 54.0\%), compared with 148,220 sequences from males (46.0\%). This female skew was strongest in the post-restriction period PR (59.0\% female) and narrowed to near-parity during spring/summer 2021 easing (47.8--51.4\% female across periods SL, L3, L21, and L0). Age composition shifted more strongly: the record was dominated overall by adults aged 25--64, but the 75+ group comprised 23.6\% of sequences in F5 and 16.9\% in T1, children aged 05--14 rose to 21.0\% in NN, and young adults aged 15--24 peaked at 35.2\% in P3.

Geographically, the record was urban-centred. Across policy periods, large and other urban areas accounted for most sequenced infections, while rural areas and small towns contributed smaller shares. At health-board level, the central-belt and north-east boards accounted for most sequences, led by Greater Glasgow and Clyde and Lothian; Orkney, Shetland, and Western Isles each contributed well under 1\% of the record. Compared with age and deprivation, the geographic profile was relatively stable across epidemic eras.

Area-level deprivation showed a modest skew toward more deprived areas. Across the full record, SIMD Q1 contributed 22.8\% of sequences compared with 18.3\% from Q5, exceeding the 20\% share expected under equal population representation. The Q1 share was highest in selected restriction and early-easing periods, including T1 (28.8\%), L2 (31.1\%), and SL (31.5\%), but was closer to the population benchmark during the Omicron response (21.9\%).

Vaccination status changed with the national rollout and is therefore summarised over rolling windows rather than as a fixed cohort attribute (Figure~\ref{fig:vaccination_context}). No sequences were linked to a recorded prior dose before the vaccination programme began in December 2020, and only 0--2.3\% were linked to a prior dose in the earliest eligible periods. This share rose thereafter, and the dose profile shifted from one-dose and two-dose records to records indicating a booster or at least three doses during the Omicron response. Among sequences from people with a recorded prior vaccination, median time since last vaccination dropped to roughly 80 days during the Omicron response and early final-easing period between November 2021 and March 2022, rose above 300 days in the third quarter of 2022, then fell to around 120 days in the final windows.

\begin{figure}[htbp]
	\centering
	\includegraphics[width=\textwidth]{fig_sequence_composition_by_policy}
	\caption[Demographic, geographic, and deprivation composition of the sequenced cohort, by epidemic era]{
		\textbf{Demographic, geographic, and deprivation composition of the sequenced cohort, by epidemic era.} Each horizontal bar shows the share of all sequences contributed by a category, stacked in chronological order by epidemic era (denoted by segment colour); the colour bar summarises the mean OxCGRT (Oxford COVID-19 Government Response Tracker) restriction stringency associated with each era. Panels: (A) sex; (B) SIMD quintile (Q1, most deprived; Q5, least deprived); (C) age group; (D) urban/rural class; and (E) health board (split across two panels for legibility).
		}\label{fig:sequence_composition_by_policy}
\end{figure}

\input{assets/scotland/tab_sequence_composition_by_policy}

\begin{figure}[htbp]
	\centering
	\includegraphics[width=\textwidth]{fig_vaccination_context}
	\caption[Vaccination profiles in the sequenced cohort]{
		\textbf{Vaccination profiles of sequenced infections across rolling 21-day windows.} (A) Window sequence counts by vaccination dose group. (B) Share of sequences from people with any recorded prior dose, one recorded dose, two recorded doses, or a recorded booster or at least three doses. (C) Median and interquartile range of days since last vaccination among sequences from people with a recorded prior vaccination (shown only for windows with at least 5 such sequences). (D) Dose-group composition over time. Shaded vertical bands denote epidemic eras. Panel backgrounds in A--C are shaded by epidemic era and coloured by the mean OxCGRT Stringency Index, as in Figure~\ref{fig:sequence_composition_by_policy} and the policy strip in Figure~\ref{fig:surveillance_timeline}A.%chktex 8
	}\label{fig:vaccination_context}
\end{figure}

Taken together, the sequenced infections formed a changing surveillance cohort rather than a fixed population sample: age, deprivation, and vaccination profiles shifted across epidemic eras, while residential geography remained urban-centred and comparatively stable. These distributions describe the sampled population used in the graph analyses below.

\subsection{Compatibility graphs show geographic structure}\label{sec:ch4_graph_assortativity}

Having characterised the sequenced cohort, we next examined how compatibility weight was distributed across population categories within graphs for each window and lineage. We summarise three features: categorical assortativity across observed attributes, variation in assortativity by window and lineage, and additional measures of graph size and weighted connectivity.

\subsubsection*{Geographic attributes dominate categorical assortativity}\label{sec:ch4_assortativity_categorical}

Categorical assortativity showed a clear spatial ordering (Table~\ref{tab:assortativity_summary}, Appendix~Figure~\ref{fig:app_assortativity_pooled_window}). The primary summary is a mean estimated by generalised least squares (GLS)\nomenclature{GLS}{generalised least squares}, accounting for the correlation between overlapping windows. It estimates average assortativity across rolling windows; window medians provide descriptive summaries of the pooled window estimates.
\input{assets/scotland/tab_assortativity_summary}

Geographic attributes had the strongest assortativity. Health board had the largest GLS mean $r$ (0.425), followed by local authority (0.305) and urban/rural class (0.189); window medians were lower but preserved the same ordering. Thus, more compatibility weight connected sequences within the same health board, local authority, or urban/rural class than expected from each category's share of total node strength. SIMD quintile, age group, and age band had positive but much smaller GLS means, while the descriptive percentile ranges across windows crossed zero for age and SIMD. Sex assortativity was approximately zero. The pooled assortativity estimates also varied across windows.

\subsubsection*{Window and lineage explain most geographic variance}\label{sec:ch4_variance_decomposition}

The variance decomposition showed a similar ordering to pooled assortativity (Table~\ref{tab:assortativity_variance_decomposition}). The primary decomposition used inverse-variance weights, so that more precisely estimated graph-level assortativity values contributed more to the summary, with weights capped at the 90th percentile to limit domination by a small number of highly precise estimates. Under this specification, the model with separate terms for surveillance window and lineage explained most variation in assortativity for geographic attributes: health board (0.857), urban/rural class (0.844), and local authority (0.837). Explained fractions were lower for SIMD quintile (0.659), age band (0.647), and age group (0.631), and lowest for sex (0.414). The same ordering was largely preserved across alternative inverse-variance weight caps (Appendix~Figure~\ref{fig:app_assortativity_variance_decomposition}). Comparisons across attributes were supported by the use of almost the same combinations of windows and lineages and similar proportions of capped weights. The larger uncertainty for sex was consistent with its weaker signal (see Appendix~Table~\ref{tab:app_assortativity_variance_decomposition_diagnostics} for definitions of the diagnostic measures).

\input{assets/scotland/tab_assortativity_variance_decomposition}

The balance between window and lineage varied by attribute. Assortativity by urban/rural class was more strongly structured by window than by lineage: the window-only model explained 0.676 of the total variation, compared with 0.552 for lineage alone. Adding lineage to the window-only model increased the fraction of total variation explained by 0.169. Health board showed the reverse pattern: the lineage-only fraction exceeded the window-only fraction (0.562 versus 0.481), and adding lineage to the window-only model increased the explained fraction by 0.377. Local authority was more balanced, while both window and lineage accounted for variation in SIMD and age assortativity. Sex was the exception, with a large unexplained fraction (0.586). After adjustment for the number of window and lineage terms, the fractions explained by the combined model were lower but preserved the same attribute ordering, indicating that the geographic pattern was not simply driven by model complexity.

Together, these results show that compatibility structure was primarily geographic and that geographic assortativity varied across windows and lineages. Assortativity by deprivation and age was weaker, while sex assortativity was close to zero.

\subsubsection*{Graph size and patterns of connections}\label{sec:ch4_topology}

Additional graph summaries describe the size and connectivity of the window-lineage graphs (Appendix~\ref{sec:app_ch4_supp_results}). Node counts, retained-edge counts, total retained compatibility weight, degree, and strength varied widely and peaked during high-volume Omicron windows (Appendix~Figure~\ref{fig:app_compatibility_topology}A--B). These measures were strongly associated with absolute sequencing volume and positive-test counts, but only weakly with sequencing coverage (Appendix~Figure~\ref{fig:app_compatibility_topology_correlations}). They therefore describe the observed sampled graphs, not the structure of the underlying transmission network.%chktex 8

Degree and strength assortativity describe how sequences are connected, rather than the size of the graph. Unweighted degree assortativity was weakly negative overall (median -0.111), whereas weighted degree assortativity and weighted strength assortativity were mostly positive (medians 0.499 and 0.206, respectively). In the weighted analyses, compatibility links tended to join sequences with similar numbers of connections or similar total connection weights. These diagnostics describe structure in the sampled compatibility graph, but they do not identify transmission hubs or individual onward transmission.

\section{Discussion}\label{sec:ch4_discussion}

This chapter examined the populations represented in the Scottish genomic surveillance record, the distribution of compatibility weight within population categories, and variation in that structure between windows and lineages. The sequenced cohort changed over time, while the largest assortativity coefficients were observed for health board, local authority, and urban/rural class. These findings connect the surveillance and population context reviewed in Chapter~\ref{ch:literature} to the pairwise relationships estimated by EpiLink.

The analysis describes additional structure among sampled infections because demographic and residential attributes were not used to calculate pairwise compatibility scores. That structure remains conditional on the EpiLink assumptions, the choice of comparison windows and lineages, and the infections included in surveillance. Its value is to establish a context for interpreting genomic clusters and to identify patterns that warrant epidemiological explanation.

\subsection{Pairwise compatibility was geographically structured}\label{sec:ch4_discussion_pairwise_epidemiology}

Health-board and local-authority assortativity showed that more compatibility weight connected sampled infections within the same administrative area than expected from the category strength shares. This is consistent with the regional circulation and lineage expansion described in Scottish and UK genomic studies \citep{daSilvaFilipe2020,duPlessis2021,Vohringer2021}. The present analysis summarises population structure in compatibility across surveillance windows; it does not reconstruct introductions or movements between areas. Geographic concentration of testing and outbreak investigation could also contribute to within-area compatibility, so the observed pattern does not distinguish regional transmission from these surveillance effects.

Assortativity by urban/rural class has a different spatial interpretation. Two people can live in the same urban/rural class while residing in distant parts of Scotland. Likewise, a shared SIMD quintile indicates similar area-level deprivation rather than a shared location. Health boards and local authorities identify areas, but residence still need not identify the place of exposure. The coefficient ranking therefore describes the chosen classifications and their category strength references; it does not rank independent effects of geography, age, sex, and deprivation on transmission.

The weak age assortativity requires consideration alongside the contact literature. Contact surveys found marked age-assortative mixing, particularly among schoolchildren and young adults \citep{Mossong2008}. Compatibility edges, however, summarise genetic and temporal evidence across all retained pairs in a window and lineage, rather than contacts recorded between individuals. Aggregation across exposure settings, the chosen age categories, and selective testing and sequencing could weaken an age pattern in these graphs. These are possible explanations requiring further investigation; the analysis cannot determine which contributes most. Weak age assortativity in compatibility therefore does not establish an absence of age structure in transmission or contacts.

The same distinction helps interpret sex and SIMD. Literature on inequalities in infection and severe outcomes motivates describing who enters surveillance \citep{PublicHealthScotland2020,Bambra2020,Niedzwiedz2020}. Such inequalities do not establish whether sampled infections with high compatibility will share the same sex or deprivation category. The cohort proportions and assortativity estimates answer complementary questions about representation and the organisation of compatibility weight. Neither alone identifies the exposure opportunities or ascertainment processes responsible for a pattern.

\subsection{Graph size and connectivity varied with surveillance conditions}\label{sec:ch4_discussion_surveillance}

Sequence volume and sequencing coverage diverged throughout the study period: shorter spring 2021 easing periods had high median coverage, while Delta and Omicron periods generated many more genomes but lower coverage relative to the larger numbers of confirmed positive tests. Here, sequencing coverage is measured against recorded positives. It does not estimate the proportion of all infections observed, and changes in testing alter the population represented by its denominator. This distinction follows the surveillance literature reviewed in Chapter~\ref{ch:literature}: diagnostic testing and sampling strategies can limit genomic surveillance even when a larger proportion of available samples is sequenced \citep{Han2023}.

Graph size and connectivity tracked sequence volume and positive-test counts more strongly than sequencing coverage. Some association with sequence volume follows from construction: more sampled genomes provide more nodes and more possible pairs. The observed correlations do not quantify separate causal effects of surveillance and transmission. They show why counts of nodes and edges, total weight, degree, and strength need to be interpreted alongside sampling volume when comparing periods.

The changing reasons for testing further qualify those comparisons. Contact tracing declined, confirmatory testing increased during parts of the Omicron response, and missing testing reasons became common after restrictions were relaxed. Age and vaccination profiles also changed, so successive graphs describe different sampled cohorts. Policy periods, testing pathways, case burden, vaccination context, and clade composition overlapped in time, and this descriptive analysis cannot separate their contributions.

\subsection{Geographic structure varied with time and lineage}\label{sec:ch4_discussion_geography_lineage}

Window and lineage together accounted for most measured variation in assortativity by health board, local authority, and urban/rural class. Assortativity by urban/rural class was more strongly associated with window, whereas health-board assortativity had a stronger lineage component. Changes in patterns by urban/rural class could accompany changing activity and sampling across epidemic phases. Differences between health boards could instead accompany the introduction and expansion of particular lineages, as suggested by previous studies of geographically structured circulation \citep{daSilvaFilipe2020,duPlessis2021}. The decomposition does not distinguish these explanations from differences in incidence, immunity, policy, or ascertainment, and its lineage terms cannot be interpreted as biological lineage effects.

Chapter~\ref{ch:literature} motivates comparisons within bounded periods and evolutionary backgrounds. The variation observed here shows why a pooled mean alone would provide an incomplete description of the sampled graphs. It does not demonstrate that the particular 21-day width, seven-day step, or Pango grouping is optimal. Those choices define the comparisons in this analysis; assessing alternatives would require sensitivity analyses of the graph summaries themselves. Leiden resolution affects the subsequent community partition and should be distinguished from choices affecting these assortativity estimates, which are calculated on the compatibility graphs before partitioning.

\subsection{Pairwise compatibility defines a baseline, not events}\label{sec:ch4_discussion_topology}

Degree and strength assortativity complement the population comparisons by describing the organisation of connections themselves. The positive weighted estimates indicate that compatibility weight tended to join genomes with similar numbers of links or similar total link weights. Such structure could accompany local lineage expansion, dense sampling of related infections, or low genetic diversity during short sampling intervals. These alternatives cannot be distinguished by connectivity alone. In particular, a genome with many compatible partners does not provide an observed count of the person's contacts or secondary infections.

Population assortativity likewise provides context for interpreting groups of compatible genomes. A geographically concentrated group may occur against a broader pattern of within-area compatibility. Establishing whether it represents a shared exposure, several transmission generations, repeated introductions, or targeted outbreak sampling requires additional evidence. The edges encode compatibility with the scenarios specified in Chapter~\ref{ch:introduction}, while subsequent cluster analyses describe how sampled groups emerge and connect through time. Neither level by itself identifies an exposure setting or realised transmission pathway.

\subsection{Strengths and limitations}\label{sec:ch4_discussion_strengths_limitations}

The main strength is the joint description of the surveillance cohort and its compatibility graphs. Linking seven data sources allowed the analysis to distinguish who was sampled from how compatibility weight was organised, while retaining sequence records, edges, and graph summaries as distinct units. Summaries within windows and lineages also made it possible to describe variation that would be obscured by a single national graph.

The analysis does not model unsequenced infections or correct for all changes in testing eligibility, sequencing priority, metadata completeness, or outbreak-focused sampling. Its findings depend on window width, Pango lineage stratification, EpiLink parameterisation, TN93 distance calculation, and sparsification. Comparisons across these choices would be needed to assess the stability of the reported graph patterns.

The uncertainty summaries also require a specific interpretation. The random-pairing reference defines assortativity, whereas the multiplier bootstrap asks how the estimate changes when the retained edge weights are perturbed. It holds the observed nodes and retained edges fixed and does not propagate uncertainty in the EpiLink parameters or the process that selected infections for sequencing. Its variances are subsequently used in pooling and precision weighting. The reported intervals are therefore conditional on the edge-perturbation model and, for pooled estimates, the pooling and working covariance assumptions. Independent edge multipliers do not by themselves account for dependence between edges sharing endpoints, and a covariance based on calendar overlap does not establish that all dependence between windows has been captured. The adequacy of these assumptions for uncertainty estimation remains a methodological question requiring further assessment.

\subsection{Implications for cluster-based analyses}\label{sec:ch4_discussion_implications}

These findings provide epidemiological context for Chapters~\ref{ch:cluster_transition_screening} and~\ref{ch:candidate_characterisation}. The association between graph size and sequencing volume motivates interpreting cluster size against its sampling context. The geographic patterns show why concentration within one health board is insufficient, by itself, to identify an unusual transmission episode. Geographic breadth likewise needs interpretation alongside lineage background, sampling density, and window timing.

Chapter~\ref{ch:cluster_transition_screening} assesses cluster size, membership novelty, and additional sampled infections in later linked clusters using its own score-specific comparison groups. The assortativity estimates reported here are not inputs to that calibration and do not constitute its null distribution. Chapter~\ref{ch:candidate_characterisation} then examines the populations and places represented among the screening outcomes. The contribution of the present chapter is to make the sampled population structure and surveillance conditions explicit before those cluster-level comparisons are interpreted.

\section{Conclusion}\label{sec:ch4_discussion_conclusions}

This chapter characterised how genetic and temporal compatibility was organised across the changing Scottish surveillance record. Compatibility weight was concentrated within health boards, local authorities, and urban/rural classes relative to their category strength references, with weaker patterns for age and deprivation and little association with sex. Surveillance window and viral lineage together accounted for most measured variation in geographic assortativity. The sequenced population changed over time, and graph size and connectivity tracked sequence volume and positive-test counts more closely than sequencing coverage.

These findings establish population and surveillance context for the subsequent cluster analyses. They motivate interpreting cluster size and residential concentration alongside sampling conditions, while leaving the exposure settings and transmission processes responsible for those patterns to further epidemiological investigation.
